% =====================================================================
% LABORATORIO LATEX: GUIDA PRATICA PER LA LEZIONE (2-3 ORE)
% Compilatore: pdfLaTeX | File immagine richiesto: immagine_esempio.png
% =====================================================================

\documentclass[a4paper,11pt]{article}
\usepackage[T1]{fontenc}
\usepackage[italian]{babel}
\usepackage[margin=2.5cm]{geometry}
\usepackage{amsmath}
\usepackage{graphicx}
\usepackage{pgfplots}
\pgfplotsset{compat=1.18}
\usepackage[colorlinks=true,linkcolor=blue,urlcolor=blue]{hyperref}

\title{Laboratorio LaTeX: elementi di un documento}
\author{Il tuo nome}
\date{\today}

% =====================================================================
% PRONTO SOCCORSO ERRORI COMUNI (Da mostrare agli studenti)
% 1. ?? nei riferimenti? -> Ricompila una seconda volta (serve a LaTeX per aggiornare i file di servizio).
% 2. Errore di sintassi improvviso? -> Controlla di non aver dimenticato una graffa chiusa '}' o un punto e virgola ';' nei grafici.
% 3. Simbolo '$' dentro l'ambiente equation o align -> Danno errore, vanno rimossi.
% 4. Nome immagine errato o maiuscole errate -> L'immagine non viene trovata.
% =====================================================================

\begin{document}

\maketitle
\tableofcontents
\clearpage

% =====================================================================
% SEZIONE 1: TESTO
% =====================================================================
\section{Testo}\label{sec:testo}
Questo è un testo normale. Una riga vuota nel codice crea un nuovo paragrafo.

Una parola in \textbf{grassetto} e una in \textit{corsivo}.

\begin{center}
Una frase centrata all'interno di un ambiente apposito.
\end{center}

Caratteri speciali protetti da backslash: 50\%, \&, \#, \$. 
La barra inversa si scrive: \textbackslash{}.

\subsection{Sottotitolo di prova}
Testo sottostante al sottotitolo.

% ESERCIZIO T1 [BASE]: Prova a scrivere due paragrafi separati da una riga vuota.
% ESERCIZIO T2 [BASE]: Inserisci una parola in grassetto e una in corsivo.


% =====================================================================
% SEZIONE 2: EQUAZIONI
% =====================================================================
\section{Equazioni}
In linea con il testo: $x^{2} + a_{1} = 5$.

Su una riga separata (senza numero):
\[
  \frac{1}{2} + \sqrt{9} = 3.5
\]

Equazione numerata automaticamente con etichetta:
\begin{equation}
  E = m c^{2}
  \label{eq:energia}
\end{equation}

Richiamo automatico dell'equazione tramite il suo rimando: l'equazione~\eqref{eq:energia} esprime l'energia.

% ESERCIZIO M1-M3 [BASE]: Prova a creare una nuova formula numerata con etichetta e richiamala nel testo.


% =====================================================================
% SEZIONE 3: ELENCHI
% =====================================================================
\section{Elenchi}
Elenco puntato (\texttt{itemize}):
\begin{itemize}
  \item Primo elemento puntato.
  \item Secondo elemento puntato.
\end{itemize}

Elenco numerato (\texttt{enumerate}):
\begin{enumerate}
  \item Primo passo logico.
  \item Secondo passo logico.
\end{enumerate}

% ESERCIZIO L1 [BASE]: Aggiungi una terza voce a entrambi gli elenchi e osserva la numerazione automatica.


% =====================================================================
% SEZIONE 4: TABELLE
% =====================================================================
\section{Tabelle}
\begin{table}[htbp]
  \centering
  \begin{tabular}{lc}
    \hline
    Voce & Quantità \\
    \hline
    Libri & 3 \\
    Quaderni & 5 \\
    \hline
  \end{tabular}
  \caption{Una tabella di esempio.}
  \label{tab:esempio}
\end{table}

% ESERCIZIO A1-A3 [BASE]: Modifica i dati della tabella e verifica che il riferimento alla tabella~\ref{tab:esempio} funzioni correttamente.


% =====================================================================
% SEZIONE 5: IMMAGINI
% =====================================================================
\section{Immagini}
\begin{figure}[htbp]
  \centering
  \includegraphics[width=0.4\textwidth]{immagine_esempio.png}
  \caption{Una descrizione chiara dell'immagine.}
  \label{fig:immagine}
\end{figure}

% ESERCIZIO I1 [BASE]: Cambia la larghezza dell'immagine da 0.4 a 0.6 e ricompila. Riferimento: la figura~\ref{fig:immagine}.


% =====================================================================
% SEZIONE 6: GRAFICI
% =====================================================================
\section{Grafici}
\begin{figure}[htbp]
\centering
\begin{tikzpicture}
\begin{axis}[
  width=0.6\textwidth,
  xlabel={$x$}, ylabel={$y$},
  domain=-2:2, samples=50]
  \addplot {2*x}; % Funzione da tracciare (ricorda il punto e virgola finale!)
\end{axis}
\end{tikzpicture}
\caption{Il grafico della funzione $y=2x$.}
\label{fig:grafico}
\end{figure}

% ESERCIZIO G1 [BASE]: Modifica la funzione in $x^2$ e aggiorna la didascalia di conseguenza.


% =====================================================================
% SEZIONE 7: RIFERIMENTI E HYPERLINK
% =====================================================================
\section{Riferimenti e Hyperlink}
Come visto nella sezione~\ref{sec:testo}, LaTeX gestisce i rimandi in automatico. 
Allo stesso modo, possiamo inserire collegamenti esterni cliccabili:
\begin{itemize}
  \item \href{https://www.overleaf.com}{Apri il sito di Overleaf}
  \item \url{https://www.latex-project.org}
\end{itemize}


% =====================================================================
% SEZIONE 8: PROVA FINALE (F1) - TEMPLATE PRONTO PER GLI STUDENTI
% =====================================================================
\clearpage
\section{Raccolta Personale (Esercitazione F1)}
% Istruzioni per gli studenti: 
% Sostituire questo testo con la propria raccolta integrando:
% - Un'equazione numerata
% - Un elenco numerato
% - Una tabella con didascalia
% - Un'immagine e un grafico
% - I rispettivi rimandi (\ref / \eqref) funzionanti e senza segnaposti '??'.

\end{document}
